\section{Titre I --- Le contrat de travail}

\article[absolue]{Forme unique}\label{art:forme-unique}
Le contrat de travail a une seule forme. Il est conclu sans détermination
de durée. Toute clause fixant un terme est réputée non écrite.
\lean{Contrat}
\theoreme{pas_de_terme}

\article[relative]{Écrit}\label{art:ecrit}
Le contrat est constaté par écrit. L'écrit indique les parties, la date de
début, la fonction, le salaire mensuel et la durée de travail convenue.
\lean{Contrat}

\article[dispositive]{Période d'essai}\label{art:essai}
Les trente premiers jours du contrat constituent une période d'essai. Un
accord d'entreprise ou le contrat peut la réduire, la supprimer ou l'allonger,
dans la limite absolue de quatre-vingt-dix jours.
\lean{Contrat.dureeEssai}
\lean{dispDureeEssai}
\theoreme{dureeEssai_le_max}
